% answers_short.tex —— 2015 年·北京卷·理科 · 简短答案速查
% 由 tools/make_short_answers.py 自 answers.tex 的 \gksolhead/\gkanswer 提取，
% 并与 2015/2015_数学真题_北京理_答案要点.md 逐题比对后生成，勿手改（改源 MD 后重跑）。

\begin{gknotice}
\textbf{使用说明}：本节只给最终答案便于快速核对，完整推导见解析版。答案为编写者独立推导，\textbf{非官方标准答案}。
\end{gknotice}

\gkpart{选择题}
\begin{center}\begin{tabular}{ccccccccc}
\toprule
题号 & 1 & 2 & 3 & 4 & 5 & 6 & 7 & 8 \\
\midrule
答案 & A & D & B & B & C & C & C & D \\
\bottomrule
\end{tabular}\end{center}

\gkpart{填空题}
\begin{itemize}[leftmargin=2.6em]
  \item[\textbf{9.}] \(40\)
  \item[\textbf{10.}] \(\dfrac{\sqrt{3}}{3}\)
  \item[\textbf{11.}] \(1\)
  \item[\textbf{12.}] \(1\)
  \item[\textbf{13.}] \(x=\dfrac{1}{2}\)，\(y=-\dfrac{1}{6}\)
  \item[\textbf{14.}] ① \(-1\)；② \(\left[\dfrac{1}{2},1\right)\cup[2,+\infty)\)
\end{itemize}

\gkpart{解答题}
\begin{itemize}[leftmargin=2.6em]
  \item[\textbf{15.}] (1) \(2\pi\)；(2) \(-1-\dfrac{\sqrt{2}}{2}\)（在 \(x=-\dfrac{3\pi}{4}\) 处取得）
  \item[\textbf{16.}] (1) \(\dfrac{3}{7}\)；(2) \(\dfrac{10}{49}\)；(3) \(a=11\) 或 \(a=18\)
  \item[\textbf{17.}] (1) 成立（\(\triangle AEF\) 为等边三角形且 \(O\) 为 \(EF\) 中点，故 \(AO\perp EF\)；由面面垂直得 \(AO\perp\) 平面 \(EFCB\)，从而 \(AO\perp BE\)）；(2) \(-\dfrac{\sqrt{5}}{5}\)（二面角为钝角，余弦值取负）；(3) \(a=\dfrac{4}{3}\)
  \item[\textbf{18.}] (1) \(y=2x\)；(2) 成立（\(f(x)-2\left(x+\dfrac{x^{3}}{3}\right)>0\) 对 \(x\in(0,1)\) 恒成立）；(3) \(k\) 的最大值为 \(2\)
  \item[\textbf{19.}] (1) \(C:\dfrac{x^{2}}{2}+y^{2}=1\)；\(M\left(\dfrac{m}{1-n},0\right)\)；(2) 存在，\(Q(0,\pm\sqrt{2})\)
  \item[\textbf{20.}] (1) \(M=\{6,12,24\}\)；(2) 成立（递推在模 \(3\) 意义下双向封闭：\(M\) 中有一个元素 \(\equiv 0\pmod 3\) 即全部 \(\equiv 0\pmod 3\)）；(3) \(8\)
\end{itemize}
